% ============================================================
% Section 5: Conclusion & Future Work
% ============================================================

\section{Conclusion}
\label{sec:conclusion}

\subsection{Summary}

We presented InCar-MSAR, the first end-to-end, open-source pipeline for
multi-speaker automatic speech recognition in automotive cabin environments.
Our three main contributions are: (1) the first published benchmark on
AISHELL-5, establishing WER and cpWER baselines; (2) empirical demonstration
that a separation-first pipeline (SepFormer + Whisper) substantially
outperforms single-channel direct transcription; and (3) a reproducible,
containerized system with a live Streamlit demo ready for investor and OEM
presentations. Experimental results on 100 sessions from AISHELL-5 eval1
confirm that each pipeline component contributes positively to overall
accuracy (see ablation in Table~\ref{tab:ablation}).

\subsection{Limitations}

Despite the encouraging results, the current system has notable limitations:

\begin{itemize}
    \item \textbf{Language scope}: The system is tuned for Mandarin Chinese
    (AISHELL-5). Extension to multilingual in-car ASR (e.g., Vietnamese,
    Japanese) would require retraining or adapting the intent keyword mapping.
    \item \textbf{Separation model mismatch}: SepFormer was pretrained on
    WSJ0-2Mix (clean, anechoic mixtures) and is applied zero-shot to
    reverberant car-cabin audio. Fine-tuning on AISHELL-5 training data
    (near-field references available) is expected to yield significant gains.
    \item \textbf{Speaker count flexibility}: The pipeline assumes $N{=}2$
    speakers by default; sessions with 3–4 simultaneous speakers degrade
    separation quality and require explicit configuration.
\end{itemize}

\subsection{Future Work}

Several promising directions emerge from this study:

\begin{enumerate}
    \item \textbf{Domain-adaptive fine-tuning}: Fine-tune SepFormer and
    Whisper jointly on AISHELL-5 near-field references (available in the
    54\,GB training split) to reduce the train-test domain gap.
    \item \textbf{On-device deployment}: Quantize the pipeline (INT8/INT4)
    for embedded deployment on automotive-grade SoCs (e.g., Qualcomm SA8295P).
    \item \textbf{LLM-based intent understanding}: Replace the keyword engine
    with a small instruction-tuned LLM (e.g., Qwen-1.5-1.8B) for richer
    intent extraction and multi-turn dialogue.
    \item \textbf{Real-time hardware integration}: Integrate with
    vehicle CAN bus simulators to enable actuator-in-the-loop evaluation.
    \item \textbf{Multilingual extension}: Leverage Whisper's multilingual
    backbone to extend the demo to Vietnamese and English cabin scenarios.
\end{enumerate}

\subsection{Broader Impact}

The open-source release of InCar-MSAR benefits the research community by
providing a reproducible baseline on AISHELL-5 and lowering the barrier to
in-car multi-speaker ASR research. Commercial applications in OEM
infotainment and ride-sharing services could improve accessibility for
hearing-impaired passengers and enhance safety by reducing driver distraction.
All third-party assets are used in accordance with their respective licenses
(AISHELL-5: CC BY-SA 4.0; Whisper: MIT; SpeechBrain: Apache 2.0).
